%% =================================================================
%% Chen et al. Measurement 264 (2026) 120300.
%% Reconstruction of page 10 (printed p. 10).
%% Contents: Fig. 13 (effect of Shore hardness on spherical cap
%% reconstruction, 8 sub-panels), Fig. 14 (effect of materials:
%% AB glue / PDMS / Ecoflex), end of Section 4.4.
%% =================================================================

\documentclass[11pt]{article}

\usepackage[margin=0.85in]{geometry}
\usepackage{amsmath, amssymb}
\usepackage{mathrsfs}
\usepackage{xcolor}
\usepackage{fancyhdr}

\pagestyle{fancy}
\fancyhf{}
\lhead{\small Z.~Chen et al.}
\rhead{\small Measurement 264 (2026) 120300}
\cfoot{\small 10 $\to$ \arabic{page}}
\renewcommand{\headrulewidth}{0.4pt}

\setcounter{page}{1}
\setcounter{equation}{1}
\setlength{\parskip}{6pt}
\setlength{\parindent}{0pt}

\begin{document}

% ---------- Fig. 13 placeholder ----------
\begin{center}
\fbox{\begin{minipage}[c][8cm][c]{0.95\textwidth}
\centering
\textbf{Fig. 13.}\ Effect of materials properties on the tactile
tomography system.\\[4pt]
\textcolor{gray}{\footnotesize
\textbf{(a)} Schematic: spherical-cap model filled with AB glues of
different Shore hardnesses (5, 10, 15, 30 HC) over an ABS
substrate.\\[3pt]
\textbf{(b)} 3D imaging of the spherical cap with a 5 HC soft
surface layer under a constant pressure of 4 MPa.\\[3pt]
\textbf{(c)} 10 HC, P = 4 MPa. \textbf{(d)} 10 HC, P = 6 MPa.\\[3pt]
\textbf{(e)} 15 HC, P = 4 MPa. \textbf{(f)} 15 HC, P = 8 MPa.\\[3pt]
\textbf{(g)} 30 HC, P = 4 MPa. \textbf{(h)} 30 HC, P = 13 MPa.\\[3pt]
Each panel renders a rainbow-colored 3D dome; Z scale 0--3.6 mm.
Harder surface layers at fixed 4 MPa give only partial cap
reconstructions; raising the probe pressure recovers the full cap.]}
\end{minipage}}
\end{center}

\vspace{0.4em}

\textbf{Fig. 13.}\ Effect of materials properties on the tactile
tomography system. (a) A spherical cap model filled with different
Shore hardness AB glues. (b) 3D imaging of the spherical cap model
with a 5 HC soft surface layer scanned by the system under a
constant pressure of 4 MPa. (c) 3D imaging of the spherical cap
model with a 10 HC soft surface layer scanned by the system under a
constant pressure of 4 MPa. (d) 3D imaging of the spherical cap
model with a 10 HC soft surface layer scanned by the system under a
constant pressure of 6 MPa. (e) 3D imaging of the spherical cap
model with a 15 HC soft surface layer scanned by the system under a
constant pressure of 4 MPa. (f) 3D imaging of the spherical cap
model with a 15 HC soft surface layer scanned by the system under a
constant pressure of 8 MPa. (g) 3D imaging of the spherical cap
model with a 30 HC soft surface layer scanned by the system under a
constant pressure of 4 MPa. (h) 3D imaging of the spherical cap
model with a 30 HC soft surface layer scanned by the system under a
constant pressure of 13 MPa.

\vspace{0.8em}

% ---------- Fig. 14 placeholder ----------
\begin{center}
\fbox{\begin{minipage}[c][5cm][c]{0.95\textwidth}
\centering
\textbf{Fig. 14.}\ Effect of materials types on the tactile
tomography system.\\[4pt]
\textcolor{gray}{\footnotesize
\textbf{(a)--(c)} Schematic cross-sections of the spherical-cap
model filled with three different soft-layer materials over the ABS
substrate: (a) AB glue, (b) PDMS, (c) Ecoflex.\\[3pt]
\textbf{(d)--(f)} 3D imaging of compression-depth distribution
when the tactile tomography system scanned the spherical cap model
with a soft surface layer of (d) AB glue, (e) PDMS, and (f)
Ecoflex. Each panel shows a clean rainbow-colored 3D dome on red
substrate (Z 0--3.6 mm), demonstrating that the system recovers the
cap geometry across all three materials.]}
\end{minipage}}
\end{center}

\vspace{0.4em}

\textbf{Fig. 14.}\ Effect of materials types on the tactile
tomography system. A spherical cap model filled with a soft surface
layer of AB glue (a), PDMS (b), and Ecoflex (c), respectively. 3D
imaging of compression depth distribution when the tactile
tomography system scanned the spherical cap model with a soft
surface layer of AB glue (d), PDMS (e), and Ecoflex (f),
respectively.

\vspace{0.6em}

% ---------- Prose tail of Section 4.4 ----------
reason for this phenomenon is that the increase in the scanning
step will result in a significant reduction in the number of
scanning points, such as when the number of scanning points is
1600, 800, and 400 when the scanning step is 0.5, 1.0, and 2.0 mm,
respectively. The low number of scanning points make the projection
data insufficient, which causes poor quality of the reconstructed
3D images. Therefore, increasing the scanning speed is one
effective strategy to achieve both high scanning efficiency and
imaging quality for the tactile tomography system.

\end{document}
